\documentclass[../../main.tex]{subfiles}
\begin{document}
\section{纤维丛}\label{section:ch10-纤维丛}

我们以简要介绍向量丛的一个重要推广来结束本章：在这种推广中，纤维允许是任意拓扑空间而非向量空间。这里我们只能触及这一课题；但纤维丛出现在流形理论的许多应用中，因此至少熟悉其定义是重要的。

设 $M$ 与 $F$ 为拓扑空间。以 $F$ 为\textbf{模型纤维}的、$M$ 上的一个\textbf{纤维丛}是一个拓扑空间 $E$，连同一个满连续映射 $\pi\colon E\to M$，满足：对每个 $x\in M$，存在 $M$ 中 $x$ 的一个邻域 $U$ 以及一个同胚 $\Phi\colon\pi^{-1}(U)\to U\times F$（称为 $E$ 在 $U$ 上的\textbf{局部平凡化}），使得下图交换：
\begin{align*}
\begin{array}{ccc}
\pi^{-1}(U)&\xrightarrow{\Phi}&U\times F\\
&\pi\searrow&\swarrow\pi_1\\
&&U.
\end{array}
\end{align*}

空间 $E$ 称为丛的全空间，$M$ 称为其底，$\pi$ 称为其投影。若 $E$、$M$、$F$ 是带边或不带边的光滑流形，$\pi$ 是光滑映射，且局部平凡化可取为微分同胚，则称为\textbf{光滑纤维丛}。

一个\textbf{平凡纤维丛}是容许在整个底空间上有局部平凡化（即整体平凡化）的丛。若它是光滑丛且整体平凡化是微分同胚，则称它\textbf{光滑平凡}。

\begin{example}[纤维丛]\label{example:ch10-纤维丛例子}
(a) 每个乘积空间 $M\times F$ 都是纤维丛，以 $\pi_1\colon M\times F\to M$ 为投影，称为\textbf{乘积纤维丛}。它有由恒等映射 $M\times F\to M\times F$ 给出的整体平凡化，故每个乘积丛都是平凡的。

(b) 每个 $k$ 秩向量丛都是以 $\mathbb{R}^k$ 为模型纤维的纤维丛。

(c) 若 $E\to S^1$ 是\cref{example:ch10-mobius丛}~中的 Möbius 丛，则在商映射 $q\colon\mathbb{R}^2\to E$ 下 $\mathbb{R}\times[-1,1]$ 的像是以 $[-1,1]$ 为模型纤维、底在 $S^1$ 上的纤维丛。它不是平凡丛。（你能证明吗？）

(d) 每个覆叠映射 $\pi\colon E\to M$ 都是模型纤维离散的纤维丛。为构造局部平凡化，令 $S$ 为与 $\pi$ 的纤维具有相同基数的离散空间。对每个被均匀覆叠的开子集 $U\subseteq M$，在 $\pi^{-1}(U)$ 的连通分支集合与 $S$ 之间选取一个双射，并令 $\Phi(x)=(\pi(x),c(x))$，其中 $c(x)$ 是 $S$ 中对应于包含 $x$ 的分支的元素。
\end{example}

随着我们继续，还会看到几个纤维丛的例子。

\section*{习题}\label{section:ch10-习题}

\begin{exercise}[Möbius 丛的光滑结构]\label{exercise:ch10-01}
设 $E$ 为\cref{example:ch10-mobius丛}~中构造的 Möbius 丛的全空间。
\begin{enumerate}[(a)]
\item 证明 $E$ 具有唯一的光滑结构，使商映射 $q\colon\mathbb{R}^2\to E$ 是光滑覆叠映射。
\item 证明 $\pi\colon E\to S^1$ 是光滑 $1$ 秩向量丛。
\item 证明它不是平凡丛。
\end{enumerate}
\end{exercise}

\begin{exercise}\label{exercise:ch10-02}
设 $E$ 为拓扑空间 $M$ 上的向量丛。证明投影映射 $\pi\colon E\to M$ 是同伦等价。
\end{exercise}

\begin{exercise}[切函子]\label{exercise:ch10-03}
设 $\mathsf{VB}$ 表示这样的范畴：其对象为光滑向量丛，态射为光滑丛同态；又设 $\mathsf{Diff}$ 表示其对象为光滑流形、态射为光滑映射的范畴。证明指派 $M\mapsto TM$、$F\mapsto dF$ 定义了一个从 $\mathsf{Diff}$ 到 $\mathsf{VB}$ 的共变函子，称为\textbf{切函子}。（用于第~303~页。）
\end{exercise}

\begin{exercise}\label{exercise:ch10-04}
补出引理~\ref{lemma:ch10-转移函数}~的证明，即证明 $\gamma\colon U\cap V\to GL(k,\mathbb{R})$ 是光滑的。（提示：用与命题~7.37~证明中相同的想法。）
\end{exercise}

\begin{exercise}[转移函数的上循环条件]\label{exercise:ch10-05}
设 $\pi\colon E\to M$ 为带边或不带边光滑流形 $M$ 上秩 $k$ 的光滑向量丛。假设 $\{U_\alpha\}_{\alpha\in A}$ 是 $M$ 的一个开覆盖，且对每个 $\alpha\in A$ 给定了 $E$ 的一个光滑局部平凡化 $\Phi_\alpha\colon\pi^{-1}(U_\alpha)\to U_\alpha\times\mathbb{R}^k$。对每个满足 $U_\alpha\cap U_\beta\neq\varnothing$ 的 $\alpha,\beta\in A$，令 $\gamma_{\alpha\beta}\colon U_\alpha\cap U_\beta\to GL(k,\mathbb{R})$ 为由\eqref{align:ch10-转移协调}~定义的转移函数。证明对所有 $\alpha,\beta,\mu\in A$ 成立
\begin{align}
\gamma_{\alpha\beta}(p)\gamma_{\beta\mu}(p)=\gamma_{\alpha\mu}(p),\qquad p\in U_\alpha\cap U_\beta\cap U_\mu.\label{align:ch10-上循环}
\end{align}
（左端的并置表示矩阵乘法。）
\end{exercise}

\begin{exercise}[向量丛构造定理]\label{exercise:ch10-06}
设 $M$ 为带边或不带边的光滑流形，$\{U_\alpha\}_{\alpha\in A}$ 为 $M$ 的一个开覆盖。假设对每个 $\alpha,\beta\in A$ 给定了光滑映射 $\gamma_{\alpha\beta}\colon U_\alpha\cap U_\beta\to GL(k,\mathbb{R})$，使得对所有 $\alpha,\beta,\mu\in A$ 满足\eqref{align:ch10-上循环}。证明存在一个光滑 $k$ 秩向量丛 $E\to M$，具有光滑局部平凡化 $\Phi_\alpha\colon\pi^{-1}(U_\alpha)\to U_\alpha\times\mathbb{R}^k$，其转移函数恰为所给的 $\gamma_{\alpha\beta}$。（提示：在 $\coprod_{\alpha\in A}(U_\alpha\times\mathbb{R}^k)$ 上定义适当的等价关系，并使用向量丛坐标图引理。）
\end{exercise}

\begin{exercise}\label{exercise:ch10-07}
计算 $T\mathbb{S}^2$ 关于由球极坐标（习题~1-7）确定的两个局部平凡化的转移函数。
\end{exercise}

\begin{exercise}\label{exercise:ch10-08}
设 $\mathsf{Vec}_1$ 为这样的范畴：其对象为有限维实向量空间，态射为线性同构。若 $F$ 是从 $\mathsf{Vec}_1$ 到自身的共变函子，则对每个有限维向量空间 $V$，我们得到一个映射 $F\colon GL(V)\to GL(F(V))$，它把每个同构 $A\colon V\to V$ 映到诱导同构 $F(A)\colon F(V)\to F(V)$。若此映射对每个 $V$ 都光滑，则称 $F$ 是\textbf{光滑函子}。给定光滑向量丛 $E\to M$ 与光滑函子 $F\colon\mathsf{Vec}_1\to\mathsf{Vec}_1$，证明存在光滑向量丛 $F(E)\to M$，其在每点 $p\in M$ 处的纤维为 $F(E_p)$。（用于第~299~页。）
\end{exercise}

\begin{exercise}[限制丛截面的扩张引理]\label{exercise:ch10-09}
设 $M$ 为光滑流形，$E\to M$ 为光滑向量丛，$S\subseteq M$ 为带边或不带边的嵌入子流形。对限制丛 $E|_S\to S$ 的任意光滑截面 $\sigma$，证明存在 $S$ 在 $M$ 中的一个邻域 $U$ 以及 $E|_U$ 的一个光滑截面 $\widetilde\sigma$，使得 $\sigma=\widetilde\sigma|_S$。若 $E$ 有正秩，证明 $E|_S$ 的每个光滑截面都能光滑地扩张到整个 $M$，当且仅当 $S$ 是真嵌入的。
\end{exercise}

\begin{exercise}\label{exercise:ch10-10}
设 $M$ 为紧光滑流形，$E\to M$ 为秩 $k$ 的光滑向量丛。利用横截性证明 $E$ 容许具有如下性质的光滑截面 $\sigma$：若 $k>\dim M$，则 $\sigma$ 处处非零；而若 $k\leqslant\dim M$，则使 $\sigma$ 为零的点集是 $M$ 的一个光滑紧余维 $k$ 子流形。用此证明 $M$ 容许只有有限多个奇点的光滑向量场。
\end{exercise}

\begin{exercise}\label{exercise:ch10-11}
证明\cref{proposition:ch10-光滑丛同态性质}：双射丛同态是丛同构。
\end{exercise}

\begin{exercise}\label{exercise:ch10-12}
设 $\pi\colon E\to M$ 与 $\widetilde\pi\colon\widetilde E\to M$ 为带边或不带边光滑流形 $M$ 上的两个光滑 $k$ 秩向量丛。假设 $\{U_\alpha\}_{\alpha\in A}$ 是 $M$ 的一个开覆盖，使得 $E$ 与 $\widetilde E$ 在每个 $U_\alpha$ 上都容许光滑局部平凡化。令 $\{\gamma_{\alpha\beta}\}$ 与 $\{\widetilde\gamma_{\alpha\beta}\}$ 分别为 $E$ 与 $\widetilde E$ 由所给局部平凡化决定的转移函数。证明 $E$ 与 $\widetilde E$ 在 $M$ 上光滑同构，当且仅当对每个 $\alpha\in A$ 存在光滑映射 $\sigma_\alpha\colon U_\alpha\to GL(k,\mathbb{R})$，使得
\begin{align*}
\widetilde\gamma_{\alpha\beta}(p)=\sigma_\alpha(p)\gamma_{\alpha\beta}(p)\sigma_\beta(p)^{-1},\qquad p\in U_\alpha\cap U_\beta.
\end{align*}
\end{exercise}

\begin{exercise}\label{exercise:ch10-13}
令 $U=\mathbb{S}^1\smallsetminus\{1\}$，$V=\mathbb{S}^1\smallsetminus\{-1\}$，并定义 $\gamma\colon U\cap V\to GL(1,\mathbb{R})$ 为
\begin{align*}
\gamma(z)=
\begin{cases}
(1),&\mathrm{Im}\,z>0,\\
(-1),&\mathrm{Im}\,z<0.
\end{cases}
\end{align*}
由习题~10-6~的结果，存在一个光滑实线丛 $F\to\mathbb{S}^1$，它在 $U$ 与 $V$ 上平凡，并以 $\gamma$ 为转移函数。证明 $F$ 在 $\mathbb{S}^1$ 上与\cref{example:ch10-mobius丛}~中的 Möbius 丛光滑同构。
\end{exercise}

\begin{exercise}\label{exercise:ch10-14}
设 $M$ 是带边或不带边的光滑流形，$S\subseteq M$ 是带边或不带边的浸入子流形。以通常方式把 $T_pS$ 等同于 $T_pM$ 的子空间，证明 $TS$ 是 $TM|_S$ 的光滑子丛。（见\cref{example:ch10-子丛例子}。）
\end{exercise}

\begin{exercise}[万有丛]\label{exercise:ch10-15}
设 $V$ 为有限维实向量空间，$G_k(V)$ 为 $V$ 的 $k$ 维子空间的 Grassmann 流形（见例~1.36）。令 $\Gamma\subseteq G_k(V)\times V$ 为子集
\begin{align*}
\Gamma=\{(S,v)\in G_k(V)\times V:v\in S\}.
\end{align*}
证明 $\Gamma$ 是乘积丛 $G_k(V)\times V\to G_k(V)$ 的一个光滑 $k$ 秩子丛，因而是 $G_k(V)$ 上的光滑 $k$ 秩向量丛。（注：$\Gamma$ 称为 $G_k(V)$ 上的\textbf{万有丛}，因为在每个点 $S\in G_k(V)$ 处的纤维恰为 $S$ 本身。）
\end{exercise}

\begin{exercise}\label{exercise:ch10-16}
证明万有丛 $G_1(\mathbb{R}^2)$ 上与 Möbius 丛光滑同构。（见习题~10-1、10-13、10-15。）
\end{exercise}

\begin{exercise}\label{exercise:ch10-17}
设 $M\subseteq\mathbb{R}^n$ 是浸入子流形。证明环境切丛 $T\mathbb{R}^n|_M$ 同构于 Whitney 和 $TM\oplus NM$，其中 $NM\to M$ 是法丛。
\end{exercise}

\begin{exercise}\label{exercise:ch10-18}
设 $S$ 是 $\mathbb{R}^n$ 的真嵌入余维 $k$ 子流形。证明下列条件等价：
\begin{enumerate}[(a)]
\item 在 $\mathbb{R}^n$ 中 $S$ 的某邻域 $U$ 上存在 $S$ 的光滑定义函数，即光滑函数 $\Phi\colon U\to\mathbb{R}^k$，使 $S$ 是 $\Phi$ 的正则水平集。
\item 法丛 $NS$ 是平凡向量丛。
\end{enumerate}
\end{exercise}

\begin{exercise}\label{exercise:ch10-19}
设 $\pi\colon E\to M$ 是以 $F$ 为纤维的纤维丛。证明下列结论：
\begin{enumerate}[(a)]
\item $\pi$ 是开商映射。
\item 若丛是光滑的，则 $\pi$ 是光滑淹没。
\item $\pi$ 是固有映射当且仅当 $F$ 是紧的。
\item $E$ 是紧的当且仅当 $M$ 与 $F$ 都紧。
\end{enumerate}
（用于第~560~页。）
\end{exercise}
\end{document}
